\chapter{Analiza i planowanie}

\section{Identyfikacja i eksploracja pomysłów}
\par Celem etapu było wygenerowanie możliwie szerokiego zbioru pomysłów dotyczących funkcjonalności i architektury projektowanego systemu, bazując na doświadczeniu członków zespołu oraz wstępnych obserwacjach problemu. Zastosowano klasyczną technikę burzy mózgów gdzie każdy z członków zespołu projektowego był zarówno pomysłodawcą jak i oceniającym. Pełniła ona nie tylko funkcję generowania pomysłów, lecz także miała istotny wymiar integracyjny, zadziałała jako mechanizm inicjujący proces formowania zespołu: stworzyła bezpieczną przestrzeń wymiany opinii, zredukowała bariery komunikacyjne oraz umożliwiła członkom zespołu na wypracowanie wspólnego sposobu dyskusji nad rozwiązaniami. Pomysły rejestrowano w ustrukturyzowanej formie (lista funkcjo, możliwe komponenty systemu, scenariusze użycia). W wyniku sesji uzystako wstępną koncepcję systemu, która była podstawą do dalszych działań.

\section{Badanie ankietowe}
\par W ramach udoskonalania konceptów oraz poszukiwania obszarów, które nie zostały w pełni pokryte przez dotychczasową analizę, przeprowadzone zostało badanie ankietowe. Szczegółowy opis wyników otrzymach z przeprowadzonej ankiety przedstawiono w załączniku "Sprawozdanie z ankiety".

\section{Planowanie pracy zespołu i środki komunikacji}
\par Na etapie planowania przyjęto iteracyjny model realizacji prac oparty na podejściu Agile, uzupełniony wybranymi praktykami SCRUM. Prace podzielono na krótkie iteracje (z małymi wyjątkami, o stałej długości: 2-tygodniowe \gls{sprint}y). Ustalono także cykliczne, cotygodniowe spotkania, w ramach których definiowano zakres zadań, monitorowano postęp oraz dokonywano przegląd rezultatów. Każde spotkanie wspierane było agendą spotkania oraz \gls{meeting minutes} w celu efektywnego przebiegu obecnych jak i przyszłych spotkań. 

\section{Narzędzia wspomagające}
\par Zespół projektowy komunikował się głównie przy użyciu komunikatora Discord na wydzielonym do tego serwerze. Do przechowywania oraz kontroli wersji kodu źródłowego systemu użyty został serwis GitHub z wydzielonym repozytroium dostępnym pod adresem url: https://github.com/Labbyn/labbyn . W celu przechowywania wszelkich wyprodukowanych podczas prac dokumentów, użyty został współdzielony Dysk Google.